\documentclass[10pt]{beamer}
\usepackage[utf8]{inputenc}
\usepackage[T1,T2A]{fontenc}
\usepackage[russian]{babel}
\usepackage{color}
\usepackage{calc}
\usepackage{graphicx}
\usepackage{epstopdf}
\usepackage{hyperref}
\hypersetup{unicode,colorlinks}
\usepackage{csquotes}
\usepackage{upquote}
\usepackage{cprotect}
\usetheme[progressbar=head,numbering=fraction,block=fill]{metropolis}
\usepackage{dejavu}
\usepackage{etoolbox}
\usepackage{bm}
\usepackage[ND]{prftree}
\usepackage[tableaux]{prooftrees}
\usepackage{mathtools} % for bigtimes
\usepackage{diagbox}
\usepackage{verbatimbox}
\usepackage{fitch}
% fitch по умолчанию печатает эти правила как $\Rightarrow$I/E; в курсе импликация всюду $\to$.
\gappto\ndrules{\def\ii{\by{$\to$I}}\def\ie{\by{$\to$E}}}
%\usepackage[defaultsans]{droidsans}
%\usepackage[defaultmono]{droidmono}
%\makeatletter
%\input droid/t1fdm.fd
%\makeatother

\makeatletter
\DeclareRobustCommand*{\bfseries}{%
    \not@math@alphabet\bfseries\mathbf
    \fontseries\bfdefault\selectfont
    \boldmath
}
\makeatother

\makeatletter
\chardef\straightquote@code=\catcode`'
\chardef\backquote@code=\catcode``
\catcode`'=\active \catcode``=\active
\patchcmd{\@noligs}
{\textasciigrave}
{\fixedtextasciigrave}
{}{}
\newcommand{\fixedtextasciigrave}{%
    \makebox[.5em]{\fontencoding{TS1}\fontfamily{fvs}\selectfont\textasciigrave}% Vera Sans
}
\catcode\lq\'=\straightquote@code
\catcode\lq\`=\backquote@code
\makeatother

\setbeamercovered{again covered={\opaqueness<1->{25}}}

\forestset{close with=\ensuremath{\bigtimes}}

\newcommand{\customframefont}[1]{
    \setbeamertemplate{itemize/enumerate body begin}{#1}
    \setbeamertemplate{itemize/enumerate subbody begin}{#1}
}

\NewEnviron{framefont}[1]{
    \customframefont{#1} % for itemize/enumerate
    {#1 % For the text outside itemize/enumerate
        \BODY
    }
    \customframefont{\normalsize}
}


\makeatletter
\DeclareRobustCommand{\rvdots}{%
    \vphantom{\int\limits^x}\smash[t]{\vdots}
}
\makeatother

% Короткое вертикальное многоточие: размеры в em, поэтому оно уменьшается вместе со шрифтом
% (у \vdots они фиксированы в pt).
\DeclareRobustCommand{\ndvdots}{%
    \mathord{\vbox{\baselineskip=.3em\lineskiplimit=0pt\kern.15em\hbox{.}\hbox{.}\hbox{.}}}}
% Подвывод в схеме правила, нарисованный как во fitch: вертикальная черта и флажок под допущением.
% Выравнивание по нижней строке, чтобы в \lor E заключения подвыводов стояли на одной линии с A \lor B.
\newcommand{\ndsubproof}[2]{{%
    \setlength{\arrayrulewidth}{.2mm}%
    \begin{array}[b]{@{}|@{\,}l@{\,}}#1\\\hline\ndvdots\\#2\end{array}}}

\newcommand{\N}{\mathbb{N}}
\newcommand{\Z}{\mathbb{Z}}
\newcommand{\Q}{\mathbb{Q}}
\newcommand{\R}{\mathbb{R}}
\newcommand{\pow}[1]{\mathcal{P}({#1})}
\newcommand{\continuum}{\mathfrak{c}}

% \reveal<спец>{...}: показать содержимое с указанного слайда, до этого оставив пустое место той же ширины.
% В отличие от \pause и \uncover не оставляет невидимых PGF-специалов; внутри выключной формулы \pause
% оставляет их в пустом абзаце после формулы, и TeX набирает лишнюю пустую строку.
% Следующий \item после такой формулы получает <.(1)-> (тот же слайд, что и \reveal<+(1)->)
% или <+(1)-> (следующий слайд); в beamer «.» "--- это значение счётчика минус один.
\newcommand<>{\reveal}[1]{\alt#2{#1}{\phantom{#1}}}

% Многобуквенные имена в формулах: курсив с нормальным кернингом (а не произведение букв).
\newcommand{\ident}[1]{\mathit{#1}}
\newcommand{\Var}{\ident{Var}}
\newcommand{\Prop}{\ident{Prop}}
\newcommand{\Const}{\ident{Const}}
\newcommand{\Fun}{\ident{Fun}}
\newcommand{\Pred}{\ident{Pred}}
\newcommand{\Term}{\ident{Term}}
\newcommand{\Frm}{\ident{Form}} % \Form занят hyperref
\newcommand{\arity}{\ident{arity}}
\newcommand{\Th}{\ident{Th}}
\newcommand{\Dom}{\ident{Dom}}
\newcommand{\Ran}{\ident{Ran}}
\newcommand{\id}{\ident{id}}
\newcommand{\Ord}{\ident{Ord}}
\newcommand{\RAA}{\ident{RAA}}
% Названия теорий и аксиом: прямой шрифт.
\newcommand{\ZF}{\mathrm{ZF}}
\newcommand{\ZFC}{\mathrm{ZFC}}
\newcommand{\AC}{\mathrm{AC}}
\newcommand{\DC}{\mathrm{DC}}
\newcommand{\CH}{\mathrm{CH}}
\newcommand{\GCH}{\mathrm{GCH}}

\author{Алексей Романов}
\subtitle{Математическая логика и теория алгоритмов}
%\logo{}
\institute{МИЭТ}
\subject{Математическая логика и теория алгоритмов}
%\setbeamercovered{transparent}
%\setbeamertemplate{navigation symbols}{}


\title{Мощности множеств}
\date{\today}

\begin{document}
\begin{frame}[plain]
    \maketitle
\end{frame}

\begin{frame}{С чем сравнить множество $X$?}
    \begin{itemize}
        \item Требуется найти $|X|$. Множество $X$ дано, а множество $Y$ для сравнения мы выбираем сами.
        \item<2-> Выбираем $Y$ известной мощности так, чтобы элементы $X$ было удобно сопоставлять с элементами $Y$.
        \item<3-> Можно перенумеровать элементы $X$ без пропусков и повторений? Тогда берём $Y=\N$ (для бесконечного списка).
        \item<4-> Для равенства мощностей достаточно биекции. А если её не видно?
    \end{itemize}
\end{frame}

\begin{frame}{Как оценить мощность сверху или снизу?}
    \begin{itemize}
        \item Как показать, что $|X|\leq|Y|$?
        \item<2-> Однозначно закодировать каждый элемент $X$ элементом $Y$: разные элементы должны иметь разные коды. Это вложение $X\to Y$.
        \item<3-> Другой путь: получить все элементы $X$ как значения функции $Y\to X$. Повторения допустимы! По критерию из лекции 6 наложение даёт верхнюю оценку (в $\ZFC$).
        \item<4-> Как показать, что $|Y|\leq|X|$? Часто достаточно найти в $X$ подмножество мощности $|Y|$. Можно и применить те же приёмы: построить вложение $Y\to X$ или наложение $X\to Y$.
        \item<5-> Если верхняя и нижняя оценки совпадают, то $|X|=|Y|$ по теореме Кантора--Бернштейна. Для этих оценок можно выбрать разные множества одной мощности.
    \end{itemize}
\end{frame}

\begin{frame}{Правило произведения}
    \begin{itemize}
        \item Выбираем $a\in A$, затем любое $b\in B$. Какова мощность множества пар $(a,b)$?
        \item<2-> По \emph{правилу произведения} из комбинаторики $|A\times B|=|A|\cdot|B|$.
        \item<3-> Сами элементы $X$ не обязаны быть парами. Если кодируем их разными парами из $A\times B$, то $|X|\leq|A|\cdot|B|$. Если использованы все пары, получаем равенство.
        \item<4-> Например, для кодов $(n,m)\in\N\times\N$ верхняя оценка равна $\aleph_0\cdot\aleph_0=\aleph_0$ (диагональная нумерация).
        \item<5-> Так же можно кодировать элементы кортежами фиксированной длины: для $r$ координат берём $r$ множителей.
        \item<6-> Произведения с бесконечным числом множителей рассмотрим позже.
    \end{itemize}
\end{frame}

\begin{frame}{Правило суммы}
    \begin{itemize}
        \item Как найти $|X|$, если известны мощности частей?
        \item<2-> Вспомним \emph{правило суммы} из комбинаторики: если $X=X_0\cup X_1$ и части не пересекаются, то $|X|=|X_0|+|X_1|$. Так же для любого конечного числа частей.
        \item<3-> Если $X$ разбито на части $X_0,X_1,\ldots$, то пишем $|X|=\sum_{n\in\N}|X_n|$.
        \item<4-> Если части не более чем счётны, но нумерации не даны, по аксиоме счётного выбора выберем по одной нумерации для каждой части.
        \item<5-> Теперь кодируем $x$ парой $(n,k)$: $n$ "--- номер части, $k$ "--- номер $x$ внутри неё. Это вложение $X\to\N\times\N$.
        \item<6-> По правилу произведения $|X|\leq\aleph_0\cdot\aleph_0=\aleph_0$. Если $X$ бесконечно, то $|X|=\aleph_0$.
    \end{itemize}
\end{frame}

\begin{frame}{Конечные двоичные слова}
    \begin{itemize}
        \item Сколько конечных слов можно составить из $0$ и $1$, если длина не ограничена?
        \item<2-> Перебираем длины, а слова одной длины выписываем в лексикографическом порядке:
        \[
        \varepsilon;\quad 0,1;\quad 00,01,10,11;\quad 000,\ldots
        \]
        \item<2-> $\varepsilon$ "--- пустое слово. Слова $1$ и $01$ различны.
        \item<3-> Слов длины $k$ ровно $2^k$. Значит, до каждой длины дойдём, и каждое слово запишем ровно один раз.
        \item<4-> Список бесконечен: в нём есть слово из $k$ нулей для каждого $k$. Мощность равна $\aleph_0$.
    \end{itemize}
\end{frame}

\begin{frame}{Конечные последовательности натуральных чисел}
    \begin{itemize}
        \only<1-6>{\item А если элементы последовательности берутся из $\N$? Можно ли опять перебирать по длине?
        \item<2-> Списки длины один не закончатся: до длины два не дойдём.
        \item<3-> Вместо длины переберём вес: $s=k+a_1+\dots+a_k$ для последовательности $(a_1,\dots,a_k)$.
        \item<4-> При $s=0$ получим пустую последовательность, при $s=1$ "--- $(0)$, при $s=2$ "--- $(1)$ и $(0,0)$.
        \item<5-> При данном $s$ длина и элементы не больше $s$. Вариантов конечное число: упорядочим их по длине, затем лексикографически.
        \item<6-> Каждый список появится ровно один раз при своём весе. Списков бесконечно много, значит, мощность $\aleph_0$.}
        \only<7->{\item<7-> Другой способ: всё же разобьём множество списков по длине. Часть $X_k$ состоит из всех списков длины $k$.
        \item<8-> $X_0$ содержит только пустой список. При $k\geq1$ множество $X_k=\N^k$ счётно: $X_1=\N$ с точностью до биекции, а каждый следующий шаг добавляет множитель $\N$.
        \item<9-> Кодируем список парой (длина, номер внутри части). В данном случае нумерацию каждой части можно построить явно, поэтому аксиома счётного выбора не нужна (дополнительное упражнение: постройте эти нумерации).
        \item<10-> Отсюда верхняя оценка $\aleph_0$. Нижнюю дают списки длины один: $n\mapsto(n)$. Значит, мощность равна $\aleph_0$.}
    \end{itemize}
\end{frame}

\begin{frame}{Сколько всего рациональных чисел?}
    \begin{itemize}
        \item Паре $(m,n)$ с $m\in\Z$ и $n\in\N_{>0}$ сопоставим дробь $m/n$. Нужно ли требовать, чтобы дробь была несократимой?
        \item<2-> Для биекции нужно устранить повторения. Для верхней оценки достаточно наложения: каждое $q\in\Q$ имеет хотя бы одну такую запись.
        \item<3-> По правилу произведения $|\Z\times\N_{>0}|=\aleph_0\cdot\aleph_0=\aleph_0$. Значит, $|\Q|\leq\aleph_0$.
        \item<4-> Нижнюю оценку даёт включение $\N\subseteq\Q$: $\aleph_0\leq|\Q|$.
        \item<5-> По теореме Кантора--Бернштейна $|\Q|=\aleph_0$.
    \end{itemize}
\end{frame}

\begin{frame}{Сколько окружностей на плоскости?}
    \begin{itemize}
        \item Какие данные определяют окружность положительного радиуса на плоскости?
        \item<2-> Центр $(u,v)$ и радиус $r$ дают код $(u,v,r)\in\R^3$.
        \item<3-> Разные окружности имеют разные коды. Для множества окружностей $X$ получаем $|X|\leq|\R^3|=\continuum$.
        \item<4-> Нужно ли теперь каждой тройке действительных чисел сопоставить свою окружность?
        \item<5-> Не обязательно: $t\in\R$ сопоставим окружность радиуса $1$ с центром $(t,0)$. Разным $t$ отвечают разные окружности.
        \item<6-> Итак, $\continuum\leq|X|\leq\continuum$.
        \item<7-> Можно и сразу получить биекцию: тройке $(u,v,r)\in\R^2\times\R_{>0}$ сопоставим окружность с центром $(u,v)$ и радиусом $r$.
    \end{itemize}
\end{frame}

\begin{frame}{Бесконечные двоичные последовательности}
    \begin{itemize}
        \item Множество всех последовательностей $b:\N\to\{0,1\}$ "--- бесконечное произведение $\prod_{n\in\N}\{0,1\}$: на каждом месте выбираем $0$ или $1$. Его мощность $2^{\aleph_0}=\continuum$.
        \item Как изменится мощность при каждом из условий?
        \begin{enumerate}
            \item Единиц в последовательности лишь конечное число.
            \item На всех чётных местах стоят нули.
        \end{enumerate}
        \item<2-> В первом случае обрежем запись после последней единицы. Если единиц нет, возьмём пустое слово.
        \item<3-> Дописав к слову нулевой хвост, восстановим последовательность. Получили вложение в счётное множество конечных слов.
        \item<4-> Для каждого $n$ есть последовательность с единственной единицей на месте $n$. Значит, мощность равна $\aleph_0$.
    \end{itemize}
\end{frame}

\begin{frame}{Нули на чётных местах}
    \begin{itemize}
        \item Второе условие: $b(2n)=0$ для всех $n\in\N$. Какие данные осталось выбрать?
        \item<2-> Значения на нечётных местах совершенно произвольны. Можно ли восстановить $b$, зная только их?
        \item<3-> Да. Любой двоичной последовательности $a$ сопоставим последовательность
        \[
           0,a(0),0,a(1),0,a(2),\ldots
        \]
        \item<4-> Обратно, из допустимой $b$ извлекаем $a(n)=b(2n+1)$. Это биекция, поэтому мощность равна $\continuum$.
    \end{itemize}
\end{frame}

\begin{frame}{Где основание, а где показатель?}
    \begin{itemize}
        \addtolength{\itemsep}{-0.6pt}
        \item Сравним $\{0,1\}\to\N$ и $\N\to\{0,1\}$. Чем задаётся функция в каждом случае?
        \item<2-> $f:\{0,1\}\to\N$ задаётся парой $(f(0),f(1))$. \uncover<3->{Любая пара подходит: $\aleph_0^2=|\N\times\N|=\aleph_0$.}
        \item<4-> $g:\N\to\{0,1\}$ задаётся бесконечной двоичной последовательностью. \uncover<5->{Здесь $2^{\aleph_0}=\continuum$.}
        \item<6-> $|A\to B|=|B|^{|A|}$: $|B|$ значений для каждого из $|A|$ аргументов. Это \emph{правило степени} из комбинаторики.
        \item<7-> А сколько бесконечных последовательностей натуральных чисел?
        \item<8-> $|\N\to\N|=\aleph_0^{\aleph_0}$. Можно ли упростить это выражение? Оценим его снизу и сверху.
        \item<8-> $\aleph_0^{\aleph_0}\geq\uncover<9->{2^{\aleph_0}=\continuum.}$
        \item<10-> С другой стороны, $\aleph_0^{\aleph_0}\leq\uncover<11->{(2^{\aleph_0})^{\aleph_0}}\uncover<12->{\overset{\text{ДЗ 6.12}}{=}2^{\aleph_0\cdot\aleph_0}}\uncover<13->{=2^{\aleph_0}.}$
        \uncover<13->{Значит, $\aleph_0^{\aleph_0}=\continuum$.}
    \end{itemize}
\end{frame}

\begin{frame}{Какие законы арифметики сохраняются?}
    \begin{itemize}
        \item Для мощностей $a,b,c$ верны, например, $a(b+c)=ab+ac$ и $(a^b)^c=a^{bc}$. Почему?
        \item<2-> Чтобы доказать каждое из этих равенств, нужно построить биекцию между множествами в левой и правой частях.
        \item<2-> Тождества с $+$ и $\cdot$, верные в $\N$, верны и для любых мощностей: после раскрытия скобок и приведения подобных слагаемых обе части совпадают.
        \item<3-> Из $a\leq b$ следуют $a+c\leq b+c$, $ac\leq bc$ и $a^c\leq b^c$. Также $c^a\leq c^b$ при $c>0$.
    \end{itemize}
\end{frame}

\begin{frame}{Строгие неравенства и сокращение}
    \begin{itemize}
        \item А если $a<b$? Будет ли $a+c<b+c$? \uncover<2->{Нет: $1<2$, но $1+\aleph_0=2+\aleph_0$. Из $a+c=b+c$ не следует $a=b$.}
        \item<3-> При умножении так же: $1\cdot\aleph_0=2\cdot\aleph_0$. Ни слагаемое, ни множитель нельзя восстановить однозначно. Поэтому для мощностей в общем случае не определяют вычитание и деление как обратные операции.
        \item<4-> Сохранится ли строгое неравенство при возведении в степень?
        \item<5-> Нет: $2^{\aleph_0}=\aleph_0^{\aleph_0}=\continuum$.
    \end{itemize}
\end{frame}

\begin{frame}{Сумма и произведение бесконечных мощностей}
    \newcommand{\cardmax}{\mathop{\scalebox{1.12}{$\max$}}\nolimits}
    \begin{itemize}
        \item Пусть хотя бы одна из мощностей $a,b$ бесконечна. В $\ZFC$ верны следующие формулы (без доказательства):
        \item $a+b=\cardmax\{a,b\}$.
        \item<2-> Если к тому же $a\ne0$ и $b\ne0$, то $a\cdot b=\cardmax\{a,b\}$.
        \item<3-> Дополнительное упражнение: считая известным, что $a\cdot a=a$ для любой бесконечной мощности $a$, докажите обе формулы.
    \end{itemize}
\end{frame}

\begin{frame}{Отрезок и интервал}
    \begin{itemize}
        \item Равномощны ли $(0,1)$ и $[0,1]$?
        \item<2-> В одну сторону есть включение $(0,1)\subseteq[0,1]$. Как построить вложение в другую сторону?
        \item<3-> Функция $x\mapsto\frac14+\frac{x}{2}$ строго возрастает и переводит $[0,1]$ в $[\frac14,\frac34]\subset(0,1)$.
        \item<4-> По теореме Кантора--Бернштейна $|(0,1)|=|[0,1]|$.
        \item<5-> А как доказать, что $|(c,d)|=|[a,b]|$ для любых $c<d$ и $a<b$?
        \item<6-> Биекция $t\mapsto c+(d-c)t$ переводит $(0,1)$ в $(c,d)$, а $t\mapsto a+(b-a)t$ переводит $[0,1]$ биективно в $[a,b]$. Значит, $|(c,d)|=|[a,b]|$.
    \end{itemize}
\end{frame}

\begin{frame}{Самостоятельная проверка}
    \begin{enumerate}
        \item Сколько конечных подмножеств у $\N$?
        \item Сколько на плоскости квадратов положительной стороны, стороны которых параллельны осям координат?
    \end{enumerate}
    \begin{itemize}
        \item Сначала выберите $Y$ и предложите биекцию или две оценки.
        \item<2-> Запишем элементы подмножества в возрастающем порядке. Это вложение в множество всех конечных списков натуральных чисел, поэтому верхняя оценка "--- $\aleph_0$. Одноэлементные подмножества дают нижнюю оценку. Ответ: $\aleph_0$.
        \item<3-> Квадрат задаётся центром и длиной стороны. Для нижней оценки достаточно единичных квадратов с центрами на одной оси. Ответ: $\continuum$.
    \end{itemize}
\end{frame}
\end{document}
